\section{Strong filter jets and the inferred memory kernel}
\label{sec:filter-memory}
We retain the finite-dimensional filter and memory consequences, with the
ambient spaces and measurability specified independently of the data.
For the homogeneous lattice, let
\[
 w_{n,\vartheta}^{\nu_n}(dz)=e^{V_n^\vartheta(z)-R_n^{\nu_n}(\vartheta)}\nu_n(dz),
 \qquad
 \Pi_{n,\vartheta}^{\nu_n}
 =(z\mapsto x_n^\vartheta+e^{t_nA_\vartheta}z)_\#w_{n,\vartheta}^{\nu_n}.
\]
Here \(V_n,R_n\) are the exact nuisance potential and its log integral
from \eqref{eq:nuisance-factorization}. All parameter differentiation
holds the observed action record fixed.

For \(j\ge0\), let \(\mathbb E_j=(C_b^{j+1}(\Hh))'\), where the test
norm is the maximum of the sup norms of all Frechet derivatives through
order \(j+1\). Let \(\mathbb E_j^0\) be the closed span in this dual of
\[
 \phi\mapsto D^\ell\phi(x)[v_1,\ldots,v_\ell],
 \quad x,v_1,\ldots,v_\ell\in\Hh,\quad0\le\ell\le j.
\]
These evaluation maps are norm continuous in all their arguments: their
variation in \(x\) is at most \(\|x-y\|\prod_k\|v_k\|\), using the
extra test derivative. Since \(\Hh\) is separable, their closed linear
span is a deterministic separable Banach space. Parameter tensors use
operator norms on symmetric multilinear maps from \(\R^5\).

\begin{lemma}[Concentrated push-forward jets]
\label{lem:strong-pushforward}
Let \(\nu\) be any probability on the radius-\(R\) Hilbert ball and
\(p_\theta\nu\) a probability with \(\theta\mapsto p_\theta\) of class
\(C^r\) in \(L^1(\nu)\). On a neighbourhood of a compact parameter set, suppose \(x_\theta\)
and \(S_\theta\) are \(C^{r+1}\), all derivatives of \(S\) through
order \(r+1\) have norm at most \(\epsilon\le1\), and set
\[
 K=1+\max_{|\alpha|\le r+1}\sup_\theta\|\partial^\alpha x_\theta\|,
 \qquad M_r=1+\max_{|\alpha|\le r}\sup_\theta
                  \|\partial^\alpha p_\theta\|_{L^1(\nu)}.
\]
For \(0\le j\le r\), the push-forward has a strong order-\(j\) derivative
in \(\mathbb E_j^0\), and
\begin{equation}
 \sup_\theta\left\|\partial_\theta^j[(x_\theta+S_\theta z)_\#(p_\theta\nu)]
                  -\partial_\theta^j\delta_{x_\theta}\right\|_{\mathbb E_j}
 \le C_r(1+R)^{r+1}K^{r+1}M_r\epsilon.
 \label{eq:strong-pushforward-bound}
\end{equation}
Under jointly measurable random dependence of the displayed derivatives,
the dual-valued fields and their parameter suprema are measurable.
\end{lemma}
\begin{proof}
For a unit test function use
\[
 \phi(x_\theta+S_\theta z)-\phi(x_\theta)
 =\int_0^1D\phi(x_\theta+tS_\theta z)[S_\theta z]dt.
\]
Every differentiated term through order \(j\) retains one factor \(S z\)
or a derivative of it; all other factors involve bounded derivatives of
\(x\), vectors of norm at most \(R\), and test derivatives through
order \(j+1\). Leibniz's rule after multiplying by \(p_\theta\) gives
\eqref{eq:strong-pushforward-bound}, linearly in the \(L^1\) bound \(M_r\).
For strong differentiability, the chain rule gives finite sums of Bochner
integrals of the evaluation maps defining \(\mathbb E_j^0\).
Uniform test-function Taylor bounds, the \(L^1\) Taylor formula for
\(p\), and dominated convergence justify the derivatives in dual norm.
Each integrand is strongly measurable in the same deterministic separable
space. Joint measurability and integration give strong measurability in
the random case. Norm continuity in \(\theta\) makes the supremum equal
to the supremum over a fixed countable dense parameter set.
\end{proof}

\begin{theorem}[Time-uniform strong filter jets]
\label{thm:filter-jets}
For every fixed \(r\) and \(q<\infty\), the homogeneous balanced experiment
satisfies, for \(0\le j\le r\),
\begin{equation}
 \sup_{E\in\mathfrak E_n}E_E\left[
 \sup_{\vartheta\in\Theta}
 \|\partial_\vartheta^j\Pi_{n,\vartheta}^{\nu_n}
       -\partial_\vartheta^j\delta_{x_n^\vartheta}\|_{\mathbb E_j}^q\right]
 \le C_{r,q}e^{-\gamma_{r,q}n}.
 \label{eq:filter-jet-decay}
\end{equation}
The claim holds for arbitrary Borel working priors on the fixed Hilbert
ball. For two such priors,
\(W_1(\Pi_{n,\vartheta}^{\nu_n^1},\Pi_{n,\vartheta}^{\nu_n^2})
 \le2MR e^{-\lambda t_n}\) deterministically.
\end{theorem}
\begin{proof}
Differentiation of a normalized exponential tilt gives centered
polynomials in derivatives of \(V_n\). For example
\(\partial w=(\partial V-w(\partial V))w\).
Induction, and the nuisance budget, therefore bound all fixed-order
\(L^1(\nu_n)\) derivatives of its density by
\(C_r(1+W_{n,r,E})^{C_r}\), uniformly in \(\vartheta\); this random
bound has every finite moment. Differentiability in total variation
follows from the same bounded Taylor formulas for the tilt. No density
of \(\nu_n\) is needed. For derivative orders above those used in the
abstract BvM, repeat the nuisance argument: the affine physical generator
has bounded derivatives of every fixed order.

Stability and bounded forcing bound every derivative of \(x_n^\vartheta\).
Derivatives of \(S_\vartheta=e^{t_nA_\vartheta}\) are at most
\(C_re^{-\lambda_rt_n}\), and \(t_n\ge n\underline\tau\) by the
common duration lower bound. Apply the push-forward lemma and then take
moments. Finitely many early \(n\) for which its \(\epsilon\le1\)
condition is not met are absorbed into \(C_{r,q}\). Coupling two
posterior initial states in the same radius-\(R\) ball proves the
Wasserstein assertion.
\end{proof}

\begin{corollary}[Current-state Gaussian image]
\label{cor:state-image}
Let \(G_n=D_\vartheta x_n^{\vartheta_0}:\R^5\to\Hh\). The joint
parameter--initial-state posterior law of
\(\sqrt n\{z(t_n)-x_n^{\vartheta_0}\}\) is at bounded-Lipschitz distance
\(o_{\mathfrak E}(1)\) from the image under \(G_n\) of
\(N(I_n^{-1}\Delta_n,I_n^{-1})\). Its comparison covariance is
\(G_n I_n^{-1}G_n^*\), of rank at most five.
\end{corollary}
\begin{proof}
On every fixed local parameter ball, Taylor's formula gives an error
\(O(|h|^2/\sqrt n)\) from replacing
\(\sqrt n(x_n^{\vartheta_0+h/\sqrt n}-x_n^{\vartheta_0})\) by \(G_nh\).
The initial-state displacement is at most \(\sqrt nMR e^{-\lambda t_n}\).
Use the parameter total-variation theorem, root-\(n\) tightness, and these
deterministic errors. The push-forward of a probability cannot increase
total variation; bounded-Lipschitz test functions control the errors.
\end{proof}

Resolve \(r=(q_0,v_0)\) and let the remaining coordinates be \(z_Q\).
Writing \(L_D\) for the half-line Dirichlet Laplacian, the unresolved
and resolved generators are
\[
 A_Q=\begin{pmatrix}0&I\\-(k_bI+\epsilon L_D)&-c_bI\end{pmatrix},
 \qquad A_P=\begin{pmatrix}0&1\\-(k+\epsilon)&-c\end{pmatrix}.
\]
The coupling \(B_Qr\) inserts \(\epsilon q_0\) into the acceleration
at site one, and \(C_Qz_Q=(0,\epsilon q_1)\). Variation of constants gives
\begin{equation}
 \dot r(t)=A_Pr(t)+\binom0{u(t)}+
 \int_0^t\mathcal K_\vartheta(t-s)r(s)ds+C_Qe^{tA_Q}z_{Q,0},
 \quad \mathcal K_\vartheta(t)=C_Qe^{tA_Q}B_Q.
 \label{eq:gle}
\end{equation}
For \(0<\omega\) smaller than a common unresolved decay rate, let
\(\mathcal X_\omega\) be the Banach space of continuous matrix kernels
with norm \(\sup_{t\ge0}e^{\omega t}\|K(t)\|\).
The map \(\mathscr M:\vartheta\mapsto\mathcal K_\vartheta\) is smooth
into this space, with bounded derivatives of every fixed order, by
Duhamel differentiation and the common exponential bound.
Its image and derivative images have deterministic separable closed
linear spans, so the random push-forwards below are Borel laws.

For large positive real \(z\), set \(p(z)=z^2+c_bz+k_b+2\epsilon\).
The tail Schur complement gives
\[
 \widehat{\mathcal K}_\vartheta(z)=\begin{pmatrix}0&0\\s(z)&0\end{pmatrix},
 \quad s(z)=\{p(z)-\sqrt{p(z)^2-4\epsilon^2}\}/2,
 \quad \widehat h_\vartheta(z)=
 \frac1{z\{z^2+cz+k+\epsilon-s(z)\}},
\]
with the square-root branch selected by \(s(z)\sim\epsilon^2/z^2\).

\begin{corollary}[Functional posterior for the memory kernel]
\label{cor:memory-posterior}
In \(\mathcal X_\omega\), the posterior law of
\(\sqrt n(\mathcal K_\vartheta-\mathcal K_{\vartheta_0})\) is at
bounded-Lipschitz distance \(o_{\mathfrak E}(1)\) from
\((D\mathscr M_{\vartheta_0})_\#N(I_n^{-1}\Delta_n,I_n^{-1})\).
The same holds for transforms on each fixed right half-plane where the
resolvents and their required derivatives are uniformly bounded.
\end{corollary}
\begin{proof}
The bounded second derivative of \(\mathscr M\) gives the uniform
Banach-valued local remainder \(O(|h|^2/\sqrt n)\).
Apply the finite-dimensional BvM and its root-\(n\) tightness.
For transforms use the resolvent identity and the same Taylor argument.
\end{proof}

\section{Finite preparation mixtures}
\label{sec:preparations}
Let a latent label \(L\in\{1,\ldots,m\}\) have prior weights \(w_j>0\),
and let the working initial-state law given \(L=j\) be \(\nu_j\) on the
fixed Hilbert ball. The parameter prior is common to all labels, and the
policy is not told the latent label. For one fixed infinite balanced
policy, summability gives a \(C^r\) almost-sure limit
\(R_\infty^{\nu_j}(\vartheta)=\lim_n R_n^{\nu_j}(\vartheta)\).

\begin{theorem}[Finite-preparation evidence mixture]
\label{thm:preparation-mixture}
The joint label and rescaled parameter posterior is at total-variation
distance \(o_P(1)\) from
\[
 \left(\sum_{j=1}^m\rho_{j,\infty}\delta_j\right)
 \otimes N(I_n^{-1}\Delta_n,I_n^{-1}),\qquad
 \rho_{j,\infty}=
 \frac{w_je^{R_\infty^{\nu_j}(\vartheta_0)}}
      {\sum_l w_le^{R_\infty^{\nu_l}(\vartheta_0)}}.
\]
For horizon-dependent policies use the finite-horizon weights with \(R_n\)
in place of \(R_\infty\); no common limiting weight across different
policies is required.
\end{theorem}
\begin{proof}
Apply the relative evidence expansion to each of the finitely many
working priors. Since
\(L_n^{\nu_j}(\vartheta_0)=L_n^0(\vartheta_0)e^{R_n^{\nu_j}(\vartheta_0)}\),
all other Laplace factors are common up to a relative \(o_P(1)\) error.
Normalization gives the weights. The nuisance potential series and its
fixed-order derivatives converge uniformly by their summable envelope;
positivity of the integral gives the asserted limit of \(R_n\).
Each conditional parameter posterior has the same Gaussian approximation.
Total variation of the finite mixture is bounded by the sum of component
errors and the \(\ell^1\) error of the weights. These weights are
finite-component Bayes factors, not thermodynamic phases.
\end{proof}
